% ==============================================================================
% SLIDES - SEMANA 04: IT BALANCED SCORECARD (BSC), PORTFÓLIO DE TI & DECISÃO
% ==============================================================================

% 1. CAPA
\senaicover
  {IT BSC \& Portfólio de TI}
  {Mapa Estratégico, Matriz do Portfólio (Weill \& Aral) e Tomada de Decisão}
  {Unidade Curricular: Governança de TI}
  {Prof. André Souza}

% 2. AGENDA
\begin{senaiframe}{Visão Geral \& Agenda}{Estrutura da Aula}
  \begin{columns}[t]
    \begin{column}{0.48\textwidth}
      \begin{tcolorbox}[senaibluecard, title={1. IT Balanced Scorecard}]
        \begin{itemize}
          \item As 4 perspectivas do BSC na TI.
          \item Relações de causa e efeito estratégicas.
        \end{itemize}
      \end{tcolorbox}
      \vspace{4pt}
      \begin{tcolorbox}[senaiambercard, title={3. Mecanismos de Decisão}]
        \begin{itemize}
          \item As 5 decisões de TI (Weill \& Ross).
          \item Alçadas formais de aprovação e comitês.
        \end{itemize}
      \end{tcolorbox}
    \end{column}
    \begin{column}{0.48\textwidth}
      \begin{tcolorbox}[senairedcard, title={2. Portfólio de TI}]
        \begin{itemize}
          \item Categorização de ativos de TI.
          \item Infraestrutura, Transacional, Informacional e Estratégico.
        \end{itemize}
      \end{tcolorbox}
      \vspace{4pt}
      \begin{tcolorbox}[senaigreencard, title={4. Capítulo 3 do PDGTI}]
        \begin{itemize}
          \item Objetivos Estratégicos de TI.
          \item Mapa Estratégico e Matriz do Portfólio.
        \end{itemize}
      \end{tcolorbox}
    \end{column}
  \end{columns}
\end{senaiframe}

% 3. ORIGEM DO BSC
\begin{senaiframe}{Teoria • Módulo 1}{Origem e Conceito do Balanced Scorecard (BSC)}
  \begin{columns}[t]
    \begin{column}{0.48\textwidth}
      \begin{tcolorbox}[senaibluecard, title={Kaplan \& Norton (1992)}]
        \begin{itemize}
          \item Superação da medição focada apenas em indicadores contábeis.
          \item Visão equilibrada do desempenho organizacional.
        \end{itemize}
      \end{tcolorbox}
    \end{column}
    \begin{column}{0.48\textwidth}
      \begin{tcolorbox}[senaicard, title={Adaptação para a TI (IT BSC)}]
        \begin{itemize}
          \item Tradução da estratégia de TI em objetivos mensuráveis.
          \item Conexão entre inovação técnica e resultado financeiro.
        \end{itemize}
      \end{tcolorbox}
    \end{column}
  \end{columns}
\end{senaiframe}

% 4. O IT BSC
\begin{senaiframe}{Teoria • Módulo 1}{As 4 Perspectivas do IT BSC}
  \begin{columns}[t]
    \begin{column}{0.48\textwidth}
      \begin{tcolorbox}[senaibluecard, title={Financeira \& Cliente}]
        \begin{itemize}
          \item \textbf{Financeira:} ROI, redução de custos e valor entregue.
          \item \textbf{Cliente / Negócio:} SLAs, valor percebido e satisfação.
        \end{itemize}
      \end{tcolorbox}
    \end{column}
    \begin{column}{0.48\textwidth}
      \begin{tcolorbox}[senaiambercard, title={Processos \& Aprendizado}]
        \begin{itemize}
          \item \textbf{Processos Internos:} Qualidade, agilidade e segurança.
          \item \textbf{Aprendizado \& Crescimento:} Capacitação da equipe.
        \end{itemize}
      \end{tcolorbox}
    \end{column}
  \end{columns}
\end{senaiframe}

% 5. PERSPECTIVA FINANCEIRA
\begin{senaiframe}{Teoria • Módulo 1}{Perspectiva 1: Financeira}
  \begin{columns}[t]
    \begin{column}{0.48\textwidth}
      \begin{tcolorbox}[senaigreencard, title={Objetivos Típicos}]
        \begin{itemize}
          \item Maximizar o retorno do capital investido em tecnologia.
          \item Otimizar os custos de operação em infraestrutura (OPEX).
        \end{itemize}
      \end{tcolorbox}
    \end{column}
    \begin{column}{0.48\textwidth}
      \begin{tcolorbox}[senaibluecard, title={Indicadores Chave}]
        \begin{itemize}
          \item Retorno sobre o Investimento (ROI de TI).
          \item Custo total de propriedade (TCO) por usuário.
        \end{itemize}
      \end{tcolorbox}
    \end{column}
  \end{columns}
\end{senaiframe}

% 6. PERSPECTIVA DO CLIENTE / NEGÓCIO
\begin{senaiframe}{Teoria • Módulo 1}{Perspectiva 2: Cliente / Negócio}
  \begin{columns}[t]
    \begin{column}{0.48\textwidth}
      \begin{tcolorbox}[senaibluecard, title={Objetivos Típicos}]
        \begin{itemize}
          \item Ser visto como parceiro estratégico das áreas de negócio.
          \item Garantir a satisfação dos usuários internos e clientes.
        \end{itemize}
      \end{tcolorbox}
    \end{column}
    \begin{column}{0.48\textwidth}
      \begin{tcolorbox}[senaicard, title={Indicadores Chave}]
        \begin{itemize}
          \item Índice de satisfação (CSAT / NPS de TI).
          \item Cumprimento dos acordos de nível de serviço (SLA).
        \end{itemize}
      \end{tcolorbox}
    \end{column}
  \end{columns}
\end{senaiframe}

% 7. PERSPECTIVA DOS PROCESSOS INTERNOS
\begin{senaiframe}{Teoria • Módulo 1}{Perspectiva 3: Processos Internos de TI}
  \begin{columns}[t]
    \begin{column}{0.48\textwidth}
      \begin{tcolorbox}[senairedcard, title={Objetivos Típicos}]
        \begin{itemize}
          \item Acelerar a entrega de software (Time-to-Market).
          \item Garantir a segurança da informação e zero bugs críticos.
        \end{itemize}
      \end{tcolorbox}
    \end{column}
    \begin{column}{0.48\textwidth}
      \begin{tcolorbox}[senaicard, title={Indicadores Chave}]
        \begin{itemize}
          \item Tempo médio de atendimento de chamados (MTTR).
          \item Taxa de defeitos em produção após o deploy.
        \end{itemize}
      \end{tcolorbox}
    \end{column}
  \end{columns}
\end{senaiframe}

% 8. PERSPECTIVA DE APRENDIZADO E CRESCIMENTO
\begin{senaiframe}{Teoria • Módulo 1}{Perspectiva 4: Aprendizado \& Crescimento}
  \begin{columns}[t]
    \begin{column}{0.48\textwidth}
      \begin{tcolorbox}[senaiambercard, title={Objetivos Típicos}]
        \begin{itemize}
          \item Capacitar a equipe técnica em tecnologias emergentes.
          \item Promover uma cultura de inovação e agilidade.
        \end{itemize}
      \end{tcolorbox}
    \end{column}
    \begin{column}{0.48\textwidth}
      \begin{tcolorbox}[senaigreencard, title={Indicadores Chave}]
        \begin{itemize}
          \item Horas de treinamento por colaborador/ano.
          \item Taxa de retenção de talentos (Turnover de TI).
        \end{itemize}
      \end{tcolorbox}
    \end{column}
  \end{columns}
\end{senaiframe}

% 9. CAUSA E EFEITO NO MAPA ESTRATÉGICO
\begin{senaiframe}{Teoria • Módulo 1}{Relações de Causa e Efeito}
  \begin{columns}[t]
    \begin{column}{0.48\textwidth}
      \begin{tcolorbox}[senaibluecard, title={Encadeamento Lógico}]
        \begin{itemize}
          \item Equipe capacitada (Aprendizado) -> Desenvolve código melhor (Processos) -> Entrega sistemas estáveis (Cliente) -> Aumenta faturamento (Financeiro).
        \end{itemize}
      \end{tcolorbox}
    \end{column}
    \begin{column}{0.48\textwidth}
      \begin{tcolorbox}[senaicard, title={Importância da Base}]
        \begin{itemize}
          \item Investir na base (capacitação) viabiliza o topo (lucro).
        \end{itemize}
      \end{tcolorbox}
    \end{column}
  \end{columns}
\end{senaiframe}

% 10. GESTÃO DE PORTFÓLIO DE TI
\begin{senaiframe}{Teoria • Módulo 2}{Gestão de Portfólio de TI}
  \begin{columns}[t]
    \begin{column}{0.48\textwidth}
      \begin{tcolorbox}[senaibluecard, title={Equilíbrio de Ativos}]
        \begin{itemize}
          \item Tratar os investimentos de TI como uma carteira de ações.
          \item Balancear riscos altos com retornos estratégicos.
        \end{itemize}
      \end{tcolorbox}
    \end{column}
    \begin{column}{0.48\textwidth}
      \begin{tcolorbox}[senaiambercard, title={Weill \& Aral}]
        \begin{itemize}
          \item Categorização em 4 tipos de investimento tecnológico.
        \end{itemize}
      \end{tcolorbox}
    \end{column}
  \end{columns}
\end{senaiframe}

% 11. AS 4 CATEGORIAS DE WEILL & ARAL
\begin{senaiframe}{Teoria • Módulo 2}{As 4 Categorias de Ativos (Weill \& Aral)}
  \begin{columns}[t]
    \begin{column}{0.48\textwidth}
      \begin{tcolorbox}[senaicard, title={Infraestrutura \& Transacional}]
        \begin{itemize}
          \item \textbf{Infraestrutura:} Sustentação (redes, servidores).
          \item \textbf{Transacional:} Automação e corte de custos (ERP).
        \end{itemize}
      \end{tcolorbox}
    \end{column}
    \begin{column}{0.48\textwidth}
      \begin{tcolorbox}[senairedcard, title={Informacional \& Estratégico}]
        \begin{itemize}
          \item \textbf{Informacional:} Inteligência e apoio à decisão (BI).
          \item \textbf{Estratégico:} Diferencial competitivo e inovação.
        \end{itemize}
      \end{tcolorbox}
    \end{column}
  \end{columns}
\end{senaiframe}

% 12. MECANISMOS DE DECISÃO
\begin{senaiframe}{Teoria • Módulo 3}{As 5 Decisões de TI (Weill \& Ross)}
  \begin{columns}[t]
    \begin{column}{0.48\textwidth}
      \begin{tcolorbox}[senaibluecard, title={Diretrizes \& Infraestrutura}]
        \begin{itemize}
          \item \textbf{1. Princípios de TI:} Papel da tecnologia na empresa.
          \item \textbf{2. Arquitetura de TI:} Padrões permitidos.
          \item \textbf{3. Infraestrutura:} Estratégia de recursos centrais.
        \end{itemize}
      \end{tcolorbox}
    \end{column}
    \begin{column}{0.48\textwidth}
      \begin{tcolorbox}[senaicard, title={Aplicações \& Investimento}]
        \begin{itemize}
          \item \textbf{4. Necessidade de Aplicações:} O que comprar/criar.
          \item \textbf{5. Investimento e Priorização:} Orçamento.
        \end{itemize}
      \end{tcolorbox}
    \end{column}
  \end{columns}
\end{senaiframe}

% 13. ESTUDO DE CASO PRÁTICO
\begin{senaiframe}{Estudo de Caso}{Mapa Estratégico de uma Fintech}
  \begin{columns}[t]
    \begin{column}{0.48\textwidth}
      \begin{tcolorbox}[senaibluecard, title={Iniciativas Mapeadas}]
        \begin{itemize}
          \item Treinar equipe em LGPD/DevSecOps (Aprendizado).
          \item Automatizar testes OWASP em código (Processos).
        \end{itemize}
      \end{tcolorbox}
    \end{column}
    \begin{column}{0.48\textwidth}
      \begin{tcolorbox}[senaigreencard, title={Resultados Alcançados}]
        \begin{itemize}
          \item API de pagamentos com 99,99\% de uptime (Cliente).
          \item Crescimento de 40\% nas transações com custo menor (Financeiro).
        \end{itemize}
      \end{tcolorbox}
    \end{column}
  \end{columns}
\end{senaiframe}

% 14. REFERÊNCIAS BIBLIOGRÁFICAS
\begin{senaiframe}{Referências Bibliográficas}{Fontes \& Leitura Recomendada}
  \begin{tcolorbox}[senaicard, title={Obras de Referência}]
    \begin{itemize}
      \item \textbf{KAPLAN, Robert S.; NORTON, David P.} *A estratégia em ação: Balanced Scorecard.* 14. ed. Rio de Janeiro: Elsevier, 1997.
      \item \textbf{WEILL, Peter; ARAL, Sinan.} *Generating Premium Returns on IT Investments.* MIT Sloan Management Review, 2006.
      \item \textbf{ITGI.} *Implementing and Nurturing IT Governance.* IT Governance Institute, 2005.
    \end{itemize}
  \end{tcolorbox}
\end{senaiframe}

% 15. APLICAÇÃO NO PDGTI
\begin{senaiframe}{Aplicação no PDGTI}{Desenvolvimento do Capítulo 3}
  \vspace{0.2cm}
  \begin{tcolorbox}[senaigreencard, title={Capítulo 3 – Objetivos Estratégicos \& Mapa Estratégico}]
    \begin{itemize}
      \item \textbf{3.1 Alinhamento de Objetivos:} Vincular Objetivos de TI ao Negócio.
      \item \textbf{3.2 Mapa Estratégico de TI:} Desenhar o mapa no modelo BSC.
      \item \textbf{3.3 Matriz do Portfólio:} Classificar projetos nas 4 categorias.
      \item \textbf{3.4 Alçadas de Decisão:} Tabela definindo responsáveis por aprovações.
    \end{itemize}
  \end{tcolorbox}
\end{senaiframe}
